\chapter{Splicing and insertion-deletion systems}
Two strings, \texttt{LABR} and \texttt{MCDN}, contain designated cut boundaries
\texttt{LA|BR} and \texttt{MC|DN}. Exchange their right-hand pieces.
The products are \texttt{LADN} and \texttt{MCBR}. Nothing was copied, yet
two new sequence arrangements appeared. Now give the first parent two
eligible cut sites. Which products are possible, and what would it mean
to say that a computation has finished?

Chapter 14 changed occupancy on a fixed backbone. Here we change the
arrangement or length of a symbolic sequence. The prerequisites are string
concatenation, contextual matching, and Chapter 12's distinction between
a language of possible species and a finite molecular inventory.
The letters in our examples are abstract symbols, not nucleotide sequences.

\section{A rule is a relation, not a laboratory instruction}
\index{splicing}\index{context}
Let $\Sigma$ be a finite alphabet and $\Sigma^*$ the finite words over it,
including the empty word $\varepsilon$. Concatenation $xy$ places $y$ after
$x$; $|x|$ counts symbols, with no implied unit of time or concentration.
Choose four context words $r=(u_1,u_2;v_1,v_2)$.
A compatible factorization of two parents is
\[
 x=a u_1u_2 b,\qquad y=c v_1v_2 d.
\]
The cuts are between $u_1$ and $u_2$, and between $v_1$ and $v_2$.
Our two-product teaching relation returns
\begin{equation}
 (x,y)\ \xRightarrow{r}\ (a u_1v_2d,\;c v_1u_2b).
 \label{eq:splice}
\end{equation}
The common formal-language convention records the first recombinant;
we explicitly retain the complementary product for inventory accounting.
The contextual cut convention follows the P\u{a}un formulation discussed
by Kari and Ng \cite{splicing}. Different splicing formalisms have different
definitions; their theorems cannot be transferred without checking hypotheses.

\Plate{01-splice}{Two symbolic parents exchange suffixes after contextual cuts.
The first product alone is a language-theoretic output; both products are
needed for the modeled material balance. Segment colors mark provenance.}{fig:splice}

Write $p=au_1$, $s=u_2b$, $q=cv_1$ and $t=v_2d$. The parents are $ps,qt$
and the products $pt,qs$. Thus
\begin{equation}
 |pt|+|qs|=|p|+|t|+|q|+|s|=|ps|+|qt|.
 \label{eq:length}
\end{equation}
A stronger invariant is the total count of each alphabet symbol across
both products. Neither invariant tells us that the reaction is feasible:
end chemistry, duplex geometry, enzyme specificity and waste are absent.
This is also not spliceosome-mediated RNA splicing.

\section{Enumerate every cut, including overlapping contexts}
\index{overlapping match}
A word of length $n$ has $n+1$ boundaries, numbered $0,\ldots,n$.
Boundary $i$ is eligible when its prefix ends in the left context and its
suffix begins with the right context. Empty contexts impose no restriction.
In \texttt{AAA}, contexts $(\texttt{A},\texttt{A})$ admit boundaries 1 and 2.
A non-overlapping pattern search can silently miss one.
\Listing{cuts}
\Listing{splice}
The return value includes both cut indices, not only product words.
Different opportunities can produce the same sequence. Deduplicating
those words is appropriate for language membership but erases multiplicity
or rate information. Even retaining every site does not determine reaction rates:
a site occurrence is not a measured binding propensity.

The reference compares all cuts of all binary words of length at most three
against an independently written boundary enumeration. It tests length
conservation for every returned pair and checks overlapping contexts.
A missing site is a completeness error even if every returned product looks plausible.

\section{A finite tube cannot reuse a consumed parent for free}
\index{inventory}\index{reaction transaction}
Let $P(w)\in\mathbb N_0$ count copies of word $w$. One selected event consumes
one copy of $x$ and one of $y$ and produces one copy of each product $z_1,z_2$:
\begin{equation}
 P'=P-\mathbf e_x-\mathbf e_y+\mathbf e_{z_1}+\mathbf e_{z_2}.
 \label{eq:inventory}
\end{equation}
Here $\mathbf e_w$ is a unit count at species $w$. If $x=y$, two copies are
required. If a product equals a parent or both products are identical,
increments still accumulate rather than overwrite one another.
\Listing{react}
The input is checked before a copied inventory is changed; failure leaves the
caller's pool intact. Three \texttt{LABR} and two \texttt{MCDN} become two
\texttt{LABR}, one \texttt{MCDN}, one \texttt{LADN} and one \texttt{MCBR}:
five copies and twenty symbols before and after.
The event is selected by cut indices. The function neither samples which
event occurs nor advances a physical clock.

Contrast this with language closure. Given axioms $L_0$ and rules $R$, define
\begin{equation}
 L_{k+1}=L_k\cup
 \{z_1,z_2:(x,y)\xRightarrow{r}(z_1,z_2),\ x,y\in L_k,\ r\in R\}.
 \label{eq:closure}
\end{equation}
Parents remain available in this mathematical definition. The union records
what can be derived, not every species simultaneously present in one tube.
A physical implementation needs replenishment or a schedule and a copy budget.
Treating one molecule as an inexhaustible reagent is not molecular parallelism.

\section{Insertion and deletion preserve context, not length}
\index{insertion}\index{deletion}
For left context $\ell$, nonempty payload $p$ and right context $r$, define
\begin{align}
 a\ell r b&\longrightarrow a\ell p r b &&\text{(insertion)},\\
 a\ell p r b&\longrightarrow a\ell r b &&\text{(deletion)}.
 \label{eq:delete}
\end{align}
The context survives; the payload appears or disappears in the word.
For \texttt{ABAB} and $(\ell,p,r)=(\texttt{A},\texttt{X},\texttt{B})$,
the one-step insertion outputs are \texttt{AXBAB} and \texttt{ABAXB}.
One step does not insert at every matching site.
\Plate{02-context}{A local insertion changes one boundary while preserving its
context. A separate donor emphasizes that physical insertion needs a source
of material; the symbolic relation contains no donor budget.}{fig:context}
\Listing{edits}
Taking all applicable sites makes this a relation with possibly many outputs.
Deleting $p$ from an inserted word can recover the original, but may also act
at another eligible site. Local inverse relations do not give a unique history.
Deleting either occurrence of \texttt{X} from \texttt{XX} with empty contexts
gives \texttt{X}; the output does not identify the removed occurrence.

Insertion raises length by $|p|$; deletion lowers it by $|p|$. An empty payload
would introduce self-transitions, so the reference rejects it. That is a scope
choice, not a universal prohibition. Physical deletion creates material outside
the retained word; a chemical balance must include released fragments.

\section{Explore reachability without confusing a cutoff with closure}
\index{reachability}\index{bounded search}
A derivation is a finite sequence of legal rewrites. Reachability asks whether
some sequence reaches a target. Begin from $\varepsilon$ and allow insertion
of either 0 or 1 at every boundary. Through depth $d$, every word of length at
most $d$ is reachable, hence
\begin{equation}
 |L_{\leq d}|=\sum_{j=0}^{d}2^j=2^{d+1}-1.
 \label{eq:growth}
\end{equation}
To prove reachability, remove one bit from a target word, construct the shorter
word inductively, then reinsert that bit. Conversely, each step adds one bit,
so no word longer than $d$ can appear in $d$ steps.
\Plate{03-branch}{Alternative insertions form a derivation graph; repeated words
merge by sequence equality. This is a graph of possibilities, not a physical
process that creates copies without reagents.}{fig:branch}

The breadth-first explorer stores a frontier and a visited set. Only the
frontier needs expansion: older words already had their successors examined.
A word-count cap raises an error rather than quietly dropping states.
At the depth limit, one extra expansion checks whether the visited set is closed.
\Listing{explore}
\begin{center}\input{\Generated/growth.tex}\end{center}
\Plate{05-budget}{Exact symbolic word counts for unrestricted binary insertion
through depth five. Exponential growth in possible words is not evidence
of an efficient molecular solver.}{fig:budget}

For a unary relation, if every frontier successor is already visited, all
visited states have successors inside the set. This proof does not justify
a binary splicing explorer: new words must also be paired with old words.
Our executable explorer is deliberately limited to unary edits.

A target found within budget has an existence witness. A target absent
from an incomplete search remains undecided. With insertions and deletions,
a short target can require a longer intermediate word; a length cap can remove
its only derivation. Report the cap alongside any negative search result.
The one-step enumeration also costs work: every boundary is examined,
and each context comparison and string construction costs more than zero.
The reference favors transparency over asymptotically optimized matching.

\section{Recover the molecular interface deliberately}
\index{cohesive end}\index{ligation}
Substring concatenation has no polarity, exposed hydroxyl or phosphate, duplex
structure, enzyme accessibility or competing ligation partner. Restore those
objects before interpreting a string program as chemistry.
\Plate{04-duplex}{Generic cohesive-end association and backbone sealing.
Dashed contacts denote pairing; solid lines denote backbone.
Association does not itself seal nicks. Enzyme-specific recognition and cut
distances are omitted; this is not a reaction protocol.}{fig:duplex}

An interface contract should identify retained fragments, complementary ends,
joinable backbone ends and recognition sites remaining after joining. Include
off-target matches, circularization and inaccessible sites as failure modes.
The same word can represent different topologies; a linear string does not
automatically model a circular plasmid.

Current Golden Gate design guidance separates Type IIS recognition sites
from fusion overhangs that order fragments and checks orientation and internal-site
compatibility \cite{neb}. This is an engineering comparison, not an implementation
of arbitrary context rewriting. Combinatorial reachability proves a statement
about chosen rules. Assembly fidelity needs a different experiment and readout.

\section{Exercises with worked answers}
\begin{enumerate}
\item \textbf{Predict.} Apply $(A,B;C,D)$ to \texttt{LABR}, \texttt{MCDN}.
\emph{Answer:} \texttt{LADN}, \texttt{MCBR}; both cut indices are 2.
\item \textbf{Prove.} Strengthen Equation~\ref{eq:length} to each symbol.
\emph{Answer:} The segments $p,s,q,t$ each occur once on each side, so their
symbol-count vectors have identical sums.
\item \textbf{Enumerate.} Count cut pairs for two \texttt{AAA} parents
with $(A,A;A,A)$. \emph{Answer:} Two choices in each parent, or four ordered pairs.
\item \textbf{Inventory.} Why cannot one \texttt{AA} copy supply both parents?
\emph{Answer:} Consumption requires two copies; subtraction would be negative.
\item \textbf{Counterexample.} Why does retaining only $z_1$ break conservation?
\emph{Answer:} Fragments in $z_2$ disappear from the model, not from reality.
\item \textbf{Implement.} Add a donor count to insertion.
\emph{Rubric:} Consume one declared donor per event, preserve caller data on failure,
account for discarded fragments and test zero donor. State remaining chemical omissions.
\item \textbf{Invert.} Is deletion injective on histories?
\emph{Answer:} No. Deleting either X from \texttt{XX} yields \texttt{X}.
\item \textbf{Derive.} Count words through depth ten.
\emph{Answer:} $2^{11}-1=2047$, including $\varepsilon$.
\item \textbf{Test.} Give a closed one-word example.
\emph{Answer:} Seed \texttt{AB} with an insertion requiring contexts X,Y has no successor.
\item \textbf{Reason.} Does absence by depth five prove non-reachability?
\emph{Answer:} Only if closure is established; otherwise longer derivations remain possible.
\item \textbf{Design.} Extend exploration to binary splicing.
\emph{Rubric:} Expand new--old, old--new and new--new pairs for asymmetric rules;
deduplicate species, preserve a derivation witness and expose budget exhaustion.
\item \textbf{Audit.} What is missing from a valid word rewrite?
\emph{Rubric:} Specify polarity, end compatibility, enzyme action, inventory,
competing reactions and a product-identity assay; distinguish these from formal proof.
\end{enumerate}
\section{Terms, reproduction and the next question}
A \Term{context} restricts a site. A \Term{recombinant} joins designated parent
segments. A \Term{closure} contains every successor of its members.
A \Term{derivation witness} is a legal path; a \Term{copy inventory} counts
available physical surrogates. README commands reproduce the generated listings
and tables. The next question reverses the viewpoint: rather than generate
words, what finite memory can recognize or transform an incoming word?
